\chapter{2003年上海卷（秋文）}
\section{填空题}
\begin{problem}
函数 \( y = \sin x \cos \left( x + \frac{\pi}{4} \right) + \cos x \sin \left( x + \frac{\pi}{4} \right) \) 的最小正周期 \( T = \)\fillinblank{}.
\end{problem}

\begin{problem}
若 \(x = \frac{\pi}{3}\) 是方程 \(2\cos (x + \alpha) = 1\) 的解，其中 \(\alpha \in (0,2\pi)\)，则 \(\alpha = \)\fillinblank{}.
\end{problem}

\begin{problem}
在等差数列 \(\{a_{n}\}\) 中，\(a_{5} = 3, a_{6} = -2\)，则 \(a_{4} + a_{5} + \cdots + a_{10} = \)\fillinblank{}.
\end{problem}

\begin{problem}
已知定点 \(A(0,1)\)，点 \(B\) 在直线 \(x + y = 0\) 上运动，当线段 \(AB\) 最短时，点 \(B\) 的坐标是\fillinblank{}.
\end{problem}

\begin{problem}
在正四棱锥 \(P - ABCD\) 中，若侧面与底面所成二面角的大小为 \(60^{\circ}\)，则异面直线 \(PA\) 与 \(BC\) 所成角的大小等于\fillinblank{}.（结果用反三角函数值表示）
\end{problem}

\begin{problem}
设集合 \(A = \{x \mid |x| < 4\}\)，\(B = \{x \mid x^2 - 4x + 3 > 0\}\)，则集合 \(\{x \mid x \in A \text{ 且 } x \notin A \cap B\} =\)\fillinblank{}.
\end{problem}

\begin{problem}
在 \(\triangle ABC\) 中，\(\sin A: \sin B: \sin C = 2:3:4\)，则 \(\angle ABC =\)\fillinblank{}.（结果用反三角函数值表示）
\end{problem}

\begin{problem}
若首项为 \(a_{1}\)，公比为 \(q\) 的等比数列 \(\{a_{n}\}\) 的前 \(n\) 项和总小于这个数列的各项和，则首项 \(a_{1}\)，公比 \(q\) 的一组取值可以是 \((a_{1}, q) =\fillinblank{}\).
\end{problem}

\begin{problem}
某国际科研合作项目成员由 11 个美国人，4 个法国人和 5 个中国人组成，现从中随机选出两位作为成果发布人，则此两人不属于同一个国家的概率为\fillinblank{}. （结果用分数表示）
\end{problem}

\begin{problem}
方程 \(x^{3} + \lg x = 18\) 的根 \(x \approx\)\fillinblank{}.（结果精确到0.1）
\end{problem}

\begin{problem}
已知点 \(A\left(0, \frac{2}{n}\right), B\left(0, -\frac{2}{n}\right), C\left(4 + \frac{2}{n}, 0\right)\)，其中 \(n\) 为正整数. 设 \(S_{n}\) 表示 \(\triangle ABC\) 外接圆的面积，则 \(\displaystyle\lim_{n \to \infty} S_{n} = \)\fillinblank{}.
\end{problem}

\begin{problem}
给出问题：\(F_{1}\)，\(F_{2}\) 是双曲线 \(\frac{x^{2}}{16} - \frac{y^{2}}{20} = 1\) 的焦点，点 \(P\) 在双曲线上. 若点 \(P\) 到焦点 \(F_{1}\) 的距离等于 9，求点 \(P\) 到焦点 \(F_{2}\) 的距离. 某学生的解答如下：双曲线的实轴长为 8，由 \(\left|PF_{1}\right| - \left|PF_{2}\right| = 8\)，即 \(\left|9 - \left|PF_{2}\right|\right| = 8\)，得 \(\left|PF_{2}\right| = 1\) 或 17. 该学生的解答是否正确？若正确，请将他的解题依据填在下面空格内，若不正确，将正确的结果填在下面空格内\fillinblank{}.
\end{problem}

\section{单选题}
\begin{problem}
下列函数中，既为偶函数又在 \((0, \pi)\) 上单调递增的是（\quad）

\choices
    {\(y = \tan |x|\)}
    {\(y = \cos (-x)\)}
    {\( y = \sin \left(x - \frac{\pi}{2}\right) \)}
    {\(y = \left|\cot \frac{x}{2}\right|\)}
\end{problem}

\begin{problem}
在下列条件中，可判断平面 \(\alpha\) 与 \(\beta\) 平行的是（\quad）

\choices
    {\(\alpha\)，\(\beta\) 都垂直于平面 \(\gamma\)}
    {\(\alpha\) 内存在不共线的三点到 \(\beta\) 的距离相等}
    {\(l, m\) 是 \(\alpha\) 内两条直线，且 \(l \parallel \beta, m \parallel \beta\)}
    {\(l, m\) 是两条异面直线，且 \(l \parallel \alpha, m \parallel \alpha, l \parallel \beta, m \parallel \beta\)}
\end{problem}

\begin{problem}
在 \(P(1,1)\)，\(Q(1,2)\)，\(M(2,3)\) 和 \(N\left(\frac{1}{2},\frac{1}{4}\right)\) 四点中，函数 \(y = a^x\) 的图象与其反函数的图象的公共点只可能是点（\quad）

\choices
    {\(P\)}
    {\(Q\)}
    {\(M\)}
    {\(N\)}
\end{problem}

\begin{problem}
\(f(x)\) 是定义在区间 \([-c, c]\) 上的奇函数，其图象如图所示：令 \(g(x) = af(x) + b\) ，则下列关于函数 \(g(x)\) 的叙述正确的是（\quad）

\bitmapfigure[page=1,pagebox=mediabox,viewport=531.0 542.5 657.0 678.5,clip,max width=0.42\linewidth,max totalheight=4.8cm]{source/普通高考/2003/2003上海文.pdf}

\choices
    {若 \(a < 0\)，则函数 \(g(x)\) 的图象关于原点对称}
    {若 \(a = -1, -2 < b < 0\)，则方程 \(g(x) = 0\) 有大于 2 的实根}
    {若 \(a \neq 0, b = 2\)，则方程 \(g(x) = 0\) 有两个实根}
    {若 \(a \geqslant 1, b < 2\)，则方程 \(g(x) = 0\) 有三个实根}
\end{problem}

\section{解答题}
\begin{problem}
已知复数 \( z_{1} = \cos \theta - \i \)，\( z_{2} = \sin \theta + \i \)，求 \( |z_{1} \cdot z_{2}| \) 的最大值和最小值.
\end{problem}

\begin{problem}
已知平行六面体 \(ABCD - A_{1}B_{1}C_{1}D_{1}\) 中，\(A_{1}A \perp\) 平面 \(ABCD, AB = 4, AD = 2\). 若 \(B_{1}D \perp BC\)，直线 \(B_{1}D\) 与平面 \(ABCD\) 所成的角等于 \(30^{\circ}\)，求平行六面体 \(ABCD - A_{1}B_{1}C_{1}D_{1}\) 的体积.

\bitmapfigure[page=1,pagebox=mediabox,viewport=998.5 650.0 1143.5 728.0,clip,max width=0.42\linewidth,max totalheight=4.8cm]{source/普通高考/2003/2003上海文.pdf}
\end{problem}

\begin{problem}
已知函数 \( f(x) = \frac{1}{x} - \log_2 \frac{1 + x}{1 - x} \)，求函数 \( f(x) \) 的定义域. 并讨论它的奇偶性和单调性.
\end{problem}

\begin{problem}
如图，某隧道设计为双向四车道，车道总宽 \(22\text{米}\)，要求通行车辆限高 \(4.5\text{米}\)，隧道全长 \(2.5\text{千米}\)，隧道的拱线近似地看成半个椭圆形状.

\begin{enumerate}
\item 若最大拱高 \(h\) 为 \(6\text{米}\)，则隧道设计的拱宽 \(l\) 是多少？
\item 若最大拱高 \(h\) 不小于 \(6\text{米}\)，则应如何设计拱高 \(h\) 和拱宽 \(l\)，才能使半个椭圆形隧道的土方工程量最小？
\end{enumerate}

【注】半个椭圆的面积公式为 \(S = \frac{\pi}{4}lh\)，柱体体积为：底面积乘以高，本题结果精确到 \(0.1\text{米}\).

\bitmapfigure[page=2,pagebox=mediabox,viewport=203.5 577.5 401.0 639.5,clip,max width=0.55\linewidth,max totalheight=4.8cm]{source/普通高考/2003/2003上海文.pdf}
\end{problem}

\begin{problem}
在以 \(O\) 为原点的直角坐标系中，点 \(A(4, -3)\) 为 \(\triangle OAB\) 的直角顶点. 已知 \(|AB| = 2|OA|\)，且点 \(B\) 的纵坐标大于零.

\begin{enumerate}
\item 求向量 \(\overrightarrow{AB}\) 的坐标；
\item 求圆 \(x^{2} - 6x + y^{2} + 2y = 0\) 关于直线 \(OB\) 对称的圆的方程；
\item 是否存在实数 \(a\)，使抛物线 \(y = ax^2 - 1\) 上总有关于直线 \(OB\) 对称的两个点？若不存在，说明理由；若存在，求 \(a\) 的取值范围.
\end{enumerate}
\end{problem}

\begin{problem}
已知数列 \(\{a_{n}\}\) （\(n\) 为正整数）是首项是 \(a_{1}\)，公比为 \(q\) 的等比数列.

\begin{enumerate}
\item 求和：\(a_{1}\mathrm C_{2}^{0} - a_{2}\mathrm C_{2}^{1} + a_{3}\mathrm C_{2}^{2}, a_{1}\mathrm C_{3}^{0} - a_{2}\mathrm C_{3}^{1} + a_{3}\mathrm C_{3}^{2} - a_{4}\mathrm C_{3}^{3}\)；
\item 由前一问的结果归纳概括出关于正整数 \(n\) 的一个结论，并加以证明；
\item 设 \(q \neq 1\)，\(S_n\) 是等比数列 \(\{a_n\}\) 的前 \(n\) 项和. 求：\(S_1\mathrm C_n^0 - S_2\mathrm C_n^1 + S_3\mathrm C_n^2 - S_4\mathrm C_n^3 + \cdots + (-1)^n S_{n+1}\mathrm C_n^n\).
\end{enumerate}
\end{problem}
